\section{Scope and manuscript contract}
\label{sec:scope}

This paper is the implementation and verification companion to the
multicomponent reacting porous-media theory paper.  Its source of truth is the
strong-form theory in that manuscript, especially the equations pulled back to
the solid skeleton \cite{foster2026reactingporous}.  The formulation inherits
the reacting-mixture setting of Drumheller and Bedford
\cite{drumheller1980thermomechanical,drumheller2000} and the finite-deformation
poroelastic structure that motivates the nonlinear Biot coefficient
\cite{foster2025}.  This document does not rederive the variational theory and
does not claim that any MOOSE object has already been implemented.  Instead, it
defines the finite-element residuals, material-property interfaces, validation
matrix, and unresolved thermodynamic derivatives needed before implementation in
MOOSE \cite{harbour2025moose}.

The computational body is the reference configuration of the solid skeleton,
\(\refdomain\).  The solid motion
\[
	\mbf x = \bs\chi_s(\mbf X,t), \qquad
	\mbf F = \Grad \bs\chi_s, \qquad
	J = \det \mbf F
\]
maps \(\refdomain\) to the current mixture domain \(\curdomain\).  All storage
terms, flux divergences, stress divergences, and source terms in this paper are
assembled per unit solid reference volume unless explicitly stated otherwise.
Spatial gradients appear only inside constitutive closures that are immediately
pulled back by \(\mbf F^{-T}\).

The first implementation target is a fully implicit finite-element solve on
\(\refdomain\) for displacement, phase/component storage variables, conversion
potentials or equivalent thermodynamic variables, phase partition variables, and
the Lagrange multipliers required by the constrained thermodynamics.  The exact
primitive variable set remains a design choice; the residual and algebraic
definitions below are written in terms of physical storage and thermodynamic
variables so that later agents can map them to MOOSE variables without changing
the mathematics or changing the reference measure of the residuals.

\subsection{Closure requirement and first closed problem}
\label{sec:closure_requirement}

For the purpose of implementation, the system is closed only when each residual
term in Section~\ref{sec:reference_solid_weak_forms} is supplied by either a
finite-element unknown, a boundary condition, or an automatic-differentiation
material property.  The general multicomponent theory is not closed by this
scaffold alone; the following list is the contract that a closed implementation
must satisfy:
\begin{itemize}[leftmargin=*]
	\item solid-skeleton kinematics: \(\mbf F\), \(J\), \(\mbf F^{-1}\), and
	the maps between spatial and reference gradients;
	\item storage variables: \(M_f^\alpha\) and \(M_s\), including the volume
	and mass-fraction constraints that define \(\phi_\xi\) and
	\(\eta_\xi^\alpha\);
	\item thermodynamic forces: \(\bar p_f\), \(\mu_\xi^\alpha\),
	\(\bar p_E\), and \(B\), with explicitly recorded held-fixed variables for
	each derivative;
	\item flux and stress closures: \(\mbf w_f\), \(\mbf W_f\),
	\(\widetilde{\bs\kappa}_f\), \(\widetilde{\mbf K}_f\),
	\(\mbf P^{\prime\prime}\), and the pressure stress contribution
	\(B\bar p_EJ\mbf F^{-T}\);
	\item conversion closures: \(\dot r_{(m)}\), \(\dot c_\xi^\alpha\), and
	any affinity or activity model used by the reaction network.
	\item phase kinematics for mobile phases: solid-frame advection of
	phase-attached scalar fields through \(\mbf W_f\).
\end{itemize}
If any item in this list is not provided, the corresponding kernel would encode
an implicit constitutive assumption and the implementation would not yet be a
closed discretization of the theory.

The first closed review problem is deliberately narrower than the full theory
but retains the full conversion structure needed for verification.  It is
isothermal, uses one mobile fluid phase and one solid skeleton phase, prescribes
a single reaction-rate law, and solves on \(\refdomain\) for displacement,
solid-reference component storage variables, phase partition variables, and the
conversion potentials required by the affinity law.  A thermodynamic material
maps those variables and \(\mbf F\) to
\(\phi_\xi\), \(\eta_\xi^\alpha\), \(\bar p_f\), \(\mu_\xi^\alpha\),
\(\bar p_E\), \(B\), and \(\dot c_\xi^\alpha\).  The first kernels therefore
include the component storage residuals, the solid storage residual, the
conversion-potential evolution equations, the reference-flux closure, and the
quasi-static skeleton momentum residual.  Additional
multifluid phase equilibrium, thermal effects, and plasticity are planned
extensions.

\subsection{Out of scope for this scaffold}
\label{sec:out_of_scope}

The following items are deliberately excluded from this initial scaffold:
\begin{itemize}[leftmargin=*]
	\item Implemented C++ kernels, materials, actions, input decks, or tests.
	\item A final choice between direct component-mass variables, pressure-composition variables, or reduced saturation-composition variables.
	\item Stabilization, upwinding, limiters, or finite-volume alternatives.
	\item Full benchmark data or production-scale reservoir examples.
\end{itemize}

Those decisions should be added only after the weak forms, thermodynamic
closures, and automatic-differentiation dependencies have been checked against
the theory manuscript.
